\section{Greedy Fundamentals \& Sorting}

\subsection{Greedy Heuristics: The Knapsack Problem}
\begin{lstlisting}
#include <iostream>
#include <vector>
#include <algorithm>

using namespace std;

struct Item {
    int value;
    int weight;
};

// Time Complexity: O(1)
bool compareByRatio(const Item& a, const Item& b) {
    double ratio_a = (double)a.value / a.weight;
    double ratio_b = (double)b.value / b.weight;
    return ratio_a > ratio_b;
}

double getMaximumKnapsackValue(vector<Item>& items, int capacity) {
    // Time Complexity: O(N log N) to sort the items
    // Space Complexity: O(log N) auxiliary space for std::sort
    sort(items.begin(), items.end(), compareByRatio);
    double total_value = 0.0;
    int current_weight = 0;
    
    // Time Complexity: O(N) - Linear scan through the sorted items
    for (const Item& item : items) {
        if (current_weight + item.weight <= capacity) {
            current_weight += item.weight;
            total_value += item.value;
        } else {
            int remaining_capacity = capacity - current_weight;
            total_value += item.value * ((double)remaining_capacity / item.weight);
            break; 
        }
    }
    return total_value;
}

int main() {
    int num_items, capacity;
    cin >> num_items >> capacity;
    // Space Complexity: O(N) where N is num_items
    vector<Item> items(num_items);
    // Time Complexity: O(N) to read items
    for (int i = 0; i < num_items; ++i) {
        cin >> items[i].value >> items[i].weight;
    }
    
    cout << getMaximumKnapsackValue(items, capacity) << "\n";
    return 0;
}
\end{lstlisting}

\subsection{Shopaholic [Kattis]}
\textbf{Concept:} Greedy Sorting. To maximize the discount, sort items descending and get the 3rd, 6th, 9th, etc., items for free.
\begin{lstlisting}
#include <iostream>
#include <vector>
#include <algorithm>

using namespace std;

long long calculateMaximumDiscount(vector<long long>& prices) {
    // Time Complexity: O(N log N)
    sort(prices.rbegin(), prices.rend()); // Sort descending
    
    long long total_discount = 0;
    // Time Complexity: O(N)
    for (size_t i = 2; i < prices.size(); i += 3) {
        total_discount += prices[i];
    }
    
    return total_discount;
}

int main() {
    int num_items;
    cin >> num_items;
    
    // Space Complexity: O(N)
    vector<long long> prices(num_items);
    for (int i = 0; i < num_items; ++i) {
        cin >> prices[i];
    }
    
    cout << calculateMaximumDiscount(prices) << "\n";
    return 0;
}
\end{lstlisting}

\subsection{Frosh Week (Inversion Counting) [Kattis]}
\textbf{Concept:} This problem asks for the minimum number of swaps to sort an array, which is exactly the number of "inversions". We use Merge Sort to count these efficiently.
\begin{lstlisting}
#include <iostream>
#include <vector>

using namespace std;
// Time Complexity: O(N log N) total for the merge sort process
// Space Complexity: O(N) for the temporary array
long long mergeAndCount(vector<int>& array, vector<int>& temp, int left, int mid, int right) {
    int i = left;
    int j = mid + 1;
    int k = left;
    long long inversions = 0;
    while (i <= mid && j <= right) {
        if (array[i] <= array[j]) {
            temp[k++] = array[i++];
        } else {
            temp[k++] = array[j++];
            // If array[i] > array[j], all remaining elements in the left half are also > array[j]
            inversions += (mid - i + 1);
        }
    }

    while (i <= mid) temp[k++] = array[i++];
    while (j <= right) temp[k++] = array[j++];
    for (i = left; i <= right; i++) array[i] = temp[i];

    return inversions;
}

long long countSwaps(vector<int>& array, vector<int>& temp, int left, int right) {
    long long inversions = 0;
    if (left < right) {
        int mid = left + (right - left) / 2;
        inversions += countSwaps(array, temp, left, mid);
        inversions += countSwaps(array, temp, mid + 1, right);
        inversions += mergeAndCount(array, temp, left, mid, right);
    }
    return inversions;
}

int main() {
    int num_students;
    if (!(cin >> num_students)) return 0;

    vector<int> student_numbers(num_students);
    vector<int> temp_array(num_students);
    for (int i = 0; i < num_students; ++i) {
        cin >> student_numbers[i];
    }

    cout << countSwaps(student_numbers, temp_array, 0, num_students - 1) << "\n";
    return 0;
}
\end{lstlisting}


\section{Interval Scheduling \& Placement}

\subsection{Greedy Scheduling (Activity Selection)}
\begin{lstlisting}
#include <iostream>
#include <vector>
#include <algorithm>

using namespace std;

struct Activity {
    int start_time;
    int end_time;
};

// Time Complexity: O(1) per comparison
bool compareActivities(Activity a, Activity b) {
    return a.end_time < b.end_time;
}

int getMaximumActivities(vector<Activity>& activity_list) {
    if (activity_list.empty()) return 0;
    // Time Complexity: O(N log N) where N is the number of activities
    // Space Complexity: O(log N) auxiliary space for std::sort (IntroSort)
    sort(activity_list.begin(), activity_list.end(), compareActivities);
    int count_selected = 1;
    int current_end_time = activity_list[0].end_time;
    
    // Time Complexity: O(N) - Single pass through the sorted activities
    for (size_t i = 1; i < activity_list.size(); ++i) {
        if (activity_list[i].start_time >= current_end_time) {
            count_selected++;
            current_end_time = activity_list[i].end_time;
        }
    }
    return count_selected;
}

int main() {
    int num_activities;
    cin >> num_activities;
    // Space Complexity: O(N) where N is num_activities
    vector<Activity> activity_list(num_activities);
    // Time Complexity: O(N) to read inputs
    for (int i = 0; i < num_activities; ++i) {
        cin >> activity_list[i].start_time >> activity_list[i].end_time;
    }
    
    cout << getMaximumActivities(activity_list) << "\n";
    return 0;
}
\end{lstlisting}

\subsection{Assigning Workstations [Kattis]}
\textbf{Concept:} Greedy Scheduling. Sort researchers by arrival time. Use a Min-Heap to track when currently used workstations will lock.
\begin{lstlisting}
#include <iostream>
#include <vector>
#include <algorithm>
#include <queue>

using namespace std;
struct Researcher {
    long long arrival;
    long long stay;
};
// Time Complexity: O(1)
bool compareArrivals(const Researcher& a, const Researcher& b) {
    return a.arrival < b.arrival;
}

int calculateSavedUnlocks(vector<Researcher>& researchers, long long lock_timeout) {
    // Time Complexity: O(N log N)
    sort(researchers.begin(), researchers.end(), compareArrivals);
    // Space Complexity: O(N) in worst case
    priority_queue<long long, vector<long long>, greater<long long>> active_machines;
    int saved_unlocks = 0;
    
    // Time Complexity: O(N log N) - N iterations, heap operations take log N
    for (const Researcher& person : researchers) {
        long long current_time = person.arrival;
        // Remove machines that have already locked
        while (!active_machines.empty() && active_machines.top() + lock_timeout < current_time) {
            active_machines.pop();
        }
        
        // If a machine is available and hasn't locked
        if (!active_machines.empty() && active_machines.top() <= current_time) {
            saved_unlocks++;
            active_machines.pop();
        }
        
        // Machine will be free after they finish
        active_machines.push(current_time + person.stay);
    }
    
    return saved_unlocks;
}

int main() {
    int num_researchers;
    long long lock_timeout;
    cin >> num_researchers >> lock_timeout;
    
    vector<Researcher> researchers(num_researchers);
    for (int i = 0; i < num_researchers; ++i) {
        cin >> researchers[i].arrival >> researchers[i].stay;
    }
    
    cout << calculateSavedUnlocks(researchers, lock_timeout) << "\n";
    return 0;
}
\end{lstlisting}

\subsection{Birds on a Wire [Kattis]}
\textbf{Concept:} Greedy Placement. Sort existing birds, then calculate how many fit in the gaps between them, including the edges.
\begin{lstlisting}
#include <iostream>
#include <vector>
#include <algorithm>

using namespace std;

int getAdditionalBirds(long long wire_length, long long distance, vector<long long>& birds) {
    if (wire_length <= 12) return 0;
    // Must be at least 6 from each pole
    
    // If no birds, calculate max possible on the whole valid wire
    if (birds.empty()) {
        return ((wire_length - 12) / distance) + 1;
    }
    
    // Time Complexity: O(N log N)
    // Space Complexity: O(log N) auxiliary
    sort(birds.begin(), birds.end());
    int new_birds = 0;
    
    // Check gap from left pole (requires 6 units)
    if (birds[0] - 6 >= distance) {
        new_birds += (birds[0] - 6) / distance;
    }
    
    // Check gaps between existing birds. Time Complexity: O(N)
    for (size_t i = 0; i < birds.size() - 1; ++i) {
        long long gap = birds[i+1] - birds[i];
        if (gap >= 2 * distance) {
            new_birds += (gap / distance) - 1;
        }
    }
    
    // Check gap to right pole (requires 6 units)
    long long last_bird = birds.back();
    if ((wire_length - 6) - last_bird >= distance) {
        new_birds += ((wire_length - 6) - last_bird) / distance;
    }
    
    return new_birds;
}

int main() {
    long long wire_length, distance;
    int num_birds;
    cin >> wire_length >> distance >> num_birds;
    
    vector<long long> birds(num_birds);
    for (int i = 0; i < num_birds; ++i) {
        cin >> birds[i];
    }
    
    cout << getAdditionalBirds(wire_length, distance, birds) << "\n";
    return 0;
}
\end{lstlisting}


\section{Priority Queues \& Data Structures}

\subsection{Huffman Codes}
\begin{lstlisting}
#include <iostream>
#include <vector>
#include <queue>

using namespace std;

struct Node {
    char character;
    int frequency;
    Node* left_child;
    Node* right_child;
};

// Time Complexity: O(1) per comparison
struct CompareNodes {
    bool operator()(Node* a, Node* b) {
        return a->frequency > b->frequency;
    }
};

Node* buildHuffmanTree(vector<char>& characters, vector<int>& frequencies) {
    // Space Complexity: O(N) where N is the number of unique characters
    priority_queue<Node*, vector<Node*>, CompareNodes> min_heap;
    // Time Complexity: O(N log N) - N insertions, each taking O(log N) time
    for (size_t i = 0; i < characters.size(); ++i) {
        // Space Complexity: O(1) per node, O(N) total across the loop
        Node* new_node = new Node{characters[i], frequencies[i], nullptr, nullptr};
        min_heap.push(new_node);
    }
    
    // Time Complexity: O(N log N) - (N-1) merges, each doing O(log N) heap operations
    while (min_heap.size() > 1) {
        Node* left = min_heap.top();
        min_heap.pop();
        Node* right = min_heap.top(); min_heap.pop();
        
        int combined_freq = left->frequency + right->frequency;
        // Space Complexity: O(1) per internal node, O(N) total for N-1 internal nodes
        Node* parent = new Node{'\0', combined_freq, left, right};
        min_heap.push(parent);
    }
    
    return min_heap.top(); 
}

int main() {
    int num_chars;
    cin >> num_chars;
    
    // Space Complexity: O(N) for characters array
    vector<char> characters(num_chars);
    // Space Complexity: O(N) for frequencies array
    vector<int> frequencies(num_chars);
    // Time Complexity: O(N) to read inputs
    for (int i = 0; i < num_chars; ++i) {
        cin >> characters[i] >> frequencies[i];
    }
    
    Node* root = buildHuffmanTree(characters, frequencies);
    return 0;
}
\end{lstlisting}

\subsection{Boiling Vegetables [Kattis]}
\textbf{Concept:} Priority Queue. Always slice the piece that is currently the largest. Keep track of the minimum piece globally to check the ratio.
\begin{lstlisting}
#include <iostream>
#include <vector>
#include <queue>

using namespace std;
struct Veggie {
    double original_weight;
    int cuts;
    double current_weight;
};
// Max-Heap comparator: order by current weight descending
struct CompareVeggie {
    bool operator()(const Veggie& a, const Veggie& b) {
        return a.current_weight < b.current_weight;
    }
};

int getMinimumCuts(double target_ratio, vector<double>& weights) {
    // Space Complexity: O(N)
    priority_queue<Veggie, vector<Veggie>, CompareVeggie> max_heap;
    double current_min_weight = 1e9;
    
    // Time Complexity: O(N log N) to build
    for (double w : weights) {
        max_heap.push({w, 1, w});
        if (w < current_min_weight) {
            current_min_weight = w;
        }
    }
    
    int total_cuts = 0;
    // Time Complexity: O(Cuts * log N). Exact bound depends on inputs.
    while (true) {
        Veggie largest = max_heap.top();
        // If the ratio between smallest and largest is valid, we are done
        if (current_min_weight / largest.current_weight > target_ratio) {
            break;
        }
        
        max_heap.pop();
        largest.cuts++;
        largest.current_weight = largest.original_weight / largest.cuts;
        
        // Update global min if the new piece is the smallest
        if (largest.current_weight < current_min_weight) {
            current_min_weight = largest.current_weight;
        }
        
        max_heap.push(largest);
        total_cuts++;
    }
    
    return total_cuts;
}

int main() {
    double target_ratio;
    int num_veggies;
    cin >> target_ratio >> num_veggies;
    
    vector<double> weights(num_veggies);
    for (int i = 0; i < num_veggies; ++i) {
        cin >> weights[i];
    }
    
    cout << getMinimumCuts(target_ratio, weights) << "\n";
    return 0;
}
\end{lstlisting}


\section{Minimum Spanning Trees (MST)}

\subsection{Prim's Algorithm (Adjacency Matrix)}
\begin{lstlisting}
#include <iostream>
#include <vector>
#include <climits>

using namespace std;
int getMinimumVertex(const vector<int>& min_edge_weights, const vector<bool>& is_in_mst, int num_vertices) {
    int min_weight = INT_MAX;
    int min_vertex = -1;
    
    // Time Complexity: O(V) - Must scan all vertices to find the minimum
    for (int v = 0; v < num_vertices; ++v) {
        if (!is_in_mst[v] && min_edge_weights[v] < min_weight) {
            min_weight = min_edge_weights[v];
            min_vertex = v;
        }
    }
    return min_vertex;
}

int calculatePrimMSTMatrix(const vector<vector<int>>& graph_matrix, int num_vertices) {
    // Space Complexity: O(V) where V is num_vertices
    vector<int> min_edge_weights(num_vertices, INT_MAX);
    // Space Complexity: O(V)
    vector<bool> is_in_mst(num_vertices, false);         
    
    int total_mst_weight = 0;
    min_edge_weights[0] = 0;
    // Time Complexity: O(V^2) total. Outer loop runs V times.
    for (int count = 0; count < num_vertices; ++count) {
        // Inner operation 1: O(V) to find the minimum vertex
        int current_u = getMinimumVertex(min_edge_weights, is_in_mst, num_vertices);
        is_in_mst[current_u] = true;
        total_mst_weight += min_edge_weights[current_u];
        
        // Inner operation 2: O(V) to update adjacent vertices
        for (int v = 0; v < num_vertices; ++v) {
            int edge_weight = graph_matrix[current_u][v];
            if (edge_weight != 0 && !is_in_mst[v] && edge_weight < min_edge_weights[v]) {
                min_edge_weights[v] = edge_weight;
            }
        }
    }
    return total_mst_weight;
}

int main() {
    int num_vertices;
    cin >> num_vertices;
    // Space Complexity: O(V^2) for the adjacency matrix
    vector<vector<int>> graph_matrix(num_vertices, vector<int>(num_vertices));
    // Time Complexity: O(V^2) to read the matrix
    for (int i = 0; i < num_vertices; ++i) {
        for (int j = 0; j < num_vertices; ++j) {
            cin >> graph_matrix[i][j];
        }
    }
    
    cout << calculatePrimMSTMatrix(graph_matrix, num_vertices) << "\n";
    return 0;
}
\end{lstlisting}

\subsection{Prim's Algorithm (Priority Queue and Adjacency List)}
\begin{lstlisting}
#include <iostream>
#include <vector>
#include <queue>

using namespace std;

struct Edge {
    int target_vertex;
    int weight;
};

// Time Complexity: O(1) per comparison
struct CompareEdge {
    bool operator()(const Edge& a, const Edge& b) {
        return a.weight > b.weight;
    }
};

int calculatePrimMSTList(const vector<vector<Edge>>& adjacency_list, int num_vertices) {
    // Space Complexity: O(E) in the worst case, where E is the number of edges
    priority_queue<Edge, vector<Edge>, CompareEdge> min_heap;
    // Space Complexity: O(V)
    vector<bool> is_in_mst(num_vertices, false);              
    
    int total_mst_weight = 0;
    int edges_added = 0;
    
    min_heap.push({0, 0});
    // Time Complexity: O(E log E) or equivalently O(E log V)
    // The while loop processes at most E edges.
    while (!min_heap.empty() && edges_added < num_vertices) {
        // Heap pop takes O(log E) time
        Edge current_edge = min_heap.top();
        min_heap.pop();
        
        int current_vertex = current_edge.target_vertex;
        
        if (is_in_mst[current_vertex]) continue;
        
        is_in_mst[current_vertex] = true;
        total_mst_weight += current_edge.weight;
        edges_added++;
        
        // Push unvisited neighbors. Over the whole algorithm, this runs E times total.
        // Heap push takes O(log E) time.
        for (const Edge& neighbor : adjacency_list[current_vertex]) {
            if (!is_in_mst[neighbor.target_vertex]) {
                min_heap.push(neighbor);
            }
        }
    }
    return total_mst_weight;
}

int main() {
    int num_vertices, num_edges;
    cin >> num_vertices >> num_edges;
    // Space Complexity: O(V + E) for the adjacency list representation
    vector<vector<Edge>> adjacency_list(num_vertices);
    // Time Complexity: O(E) to read edges
    for (int i = 0; i < num_edges; ++i) {
        int u, v, weight;
        cin >> u >> v >> weight;
        adjacency_list[u].push_back({v, weight});
        adjacency_list[v].push_back({u, weight});
    }
    
    cout << calculatePrimMSTList(adjacency_list, num_vertices) << "\n";
    return 0;
}
\end{lstlisting}

\subsection{Kruskal's Algorithm}
\subsubsection{Bases for Kruskal's Algorithm}
\begin{itemize}
    \item \textbf{The Strategy:} Sort all edges in non-decreasing order of their weight. Pick the smallest edge. If it forms a cycle with the spanning tree formed so far, discard it. If not, include it.
    \item \textbf{Cycle Detection:} Uses the Disjoint Set Union (DSU) data structure. If both vertices of an edge belong to the same set, adding the edge will create a cycle.
    \item \textbf{Complexity:} $O(E \log E)$ or $O(E \log V)$. It is generally preferred for sparse graphs where edges can be easily sorted.
\end{itemize}

\subsubsection{Implementation: Union-Find and Edge Sorting}
\begin{lstlisting}
#include <iostream>
#include <vector>
#include <algorithm>

using namespace std;

struct Edge {
    int source;
    int destination;
    int weight;
};

// Time Complexity: O(1)
bool compareEdges(const Edge& a, const Edge& b) {
    return a.weight < b.weight;
}

// --- Disjoint Set Union (DSU) Helper Functions ---

// Time Complexity: O(alpha(V)) ~ O(1) due to path compression
int findParent(int vertex, vector<int>& parent_array) {
    if (parent_array[vertex] == vertex) {
        return vertex;
    }
    // Path compression: attach the node directly to the root
    return parent_array[vertex] = findParent(parent_array[vertex], parent_array);
}

// Time Complexity: O(alpha(V)) ~ O(1)
void unionSets(int vertex_u, int vertex_v, vector<int>& parent_array, vector<int>& rank_array) {
    int root_u = findParent(vertex_u, parent_array);
    int root_v = findParent(vertex_v, parent_array);
    
    if (root_u != root_v) {
        // Union by rank: attach smaller tree under root of deeper tree
        if (rank_array[root_u] < rank_array[root_v]) {
            parent_array[root_u] = root_v;
        } else if (rank_array[root_u] > rank_array[root_v]) {
            parent_array[root_v] = root_u;
        } else {
            parent_array[root_v] = root_u;
            rank_array[root_u]++;
        }
    }
}

// --- Main Kruskal Logic ---

int calculateKruskalMST(vector<Edge>& edge_list, int num_vertices) {
    // Time Complexity: O(E log E) to sort the edges
    // Space Complexity: O(log E) auxiliary space for std::sort
    sort(edge_list.begin(), edge_list.end(), compareEdges);
    // Space Complexity: O(V) for parent and rank arrays
    vector<int> parent_array(num_vertices);
    vector<int> rank_array(num_vertices, 0);
    // Initialize DSU: Every vertex is its own parent initially
    for (int i = 0; i < num_vertices; ++i) {
        parent_array[i] = i;
    }
    
    int total_mst_weight = 0;
    int edges_added = 0;
    // Time Complexity: O(E * alpha(V)) ~ O(E) for the loop
    for (const Edge& edge : edge_list) {
        // Stop early if we've added V-1 edges (a complete spanning tree)
        if (edges_added == num_vertices - 1) break;
        int root_source = findParent(edge.source, parent_array);
        int root_dest = findParent(edge.destination, parent_array);
        // If they don't share a parent, adding this edge won't create a cycle
        if (root_source != root_dest) {
            total_mst_weight += edge.weight;
            unionSets(root_source, root_dest, parent_array, rank_array);
            edges_added++;
        }
    }
    
    return total_mst_weight;
}

int main() {
    int num_vertices, num_edges;
    cin >> num_vertices >> num_edges;
    // Space Complexity: O(E) to store the list of all edges
    vector<Edge> edge_list(num_edges);
    // Time Complexity: O(E) to read edges
    for (int i = 0; i < num_edges; ++i) {
        cin >> edge_list[i].source >> edge_list[i].destination >> edge_list[i].weight;
    }
    
    cout << calculateKruskalMST(edge_list, num_vertices) << "\n";
    return 0;
}
\end{lstlisting}

\subsection{Lost Map [Kattis]}
\textbf{Concept:} Kruskal's MST from an Adjacency Matrix. The shortest paths between adjacent villages form the MST of the complete graph.
\begin{lstlisting}
#include <iostream>
#include <vector>
#include <algorithm>

using namespace std;
struct Edge {
    int u, v, weight;
};
bool compareEdges(const Edge& a, const Edge& b) {
    return a.weight < b.weight;
}

// Time Complexity: O(1) amortized
int findParent(int i, vector<int>& parent) {
    if (parent[i] == i) return i;
    return parent[i] = findParent(parent[i], parent);
}

void printLostMapEdges(vector<Edge>& edges, int num_vertices) {
    // Time Complexity: O(V^2 log(V^2)) to sort edges
    sort(edges.begin(), edges.end(), compareEdges);
    // Space Complexity: O(V)
    vector<int> parent(num_vertices + 1);
    for (int i = 1; i <= num_vertices; ++i) parent[i] = i;
    
    int edges_printed = 0;
    // Time Complexity: O(V^2) loop iterations
    for (const Edge& e : edges) {
        if (edges_printed == num_vertices - 1) break;
        int root_u = findParent(e.u, parent);
        int root_v = findParent(e.v, parent);
        if (root_u != root_v) {
            cout << e.u << " " << e.v << "\n";
            parent[root_u] = root_v;
            edges_printed++;
        }
    }
}

int main() {
    int num_vertices;
    cin >> num_vertices;
    // Space Complexity: O(V^2)
    vector<Edge> edges;
    for (int i = 1; i <= num_vertices; ++i) {
        for (int j = 1; j <= num_vertices; ++j) {
            int weight;
            cin >> weight;
            // Only add upper triangle of the matrix to avoid duplicates
            if (i < j) {
                edges.push_back({i, j, weight});
            }
        }
    }
    
    printLostMapEdges(edges, num_vertices);
    return 0;
}
\end{lstlisting}


\section{Ad-Hoc / Mathematical Greedy}

\subsection{Ants [Kattis]}
\textbf{Concept:} Independent movement. When ants collide, they bounce, which is mathematically identical to them passing through each other.
Min time: The last ant to fall off if they all walk to their *closest* edge.
Max time: The last ant to fall off if they all walk to their *furthest* edge.
\begin{lstlisting}
#include <iostream>
#include <vector>
#include <algorithm>

using namespace std;

void calculateAntTimes(int pole_length, const vector<int>& ant_positions) {
    int earliest_time = 0;
    int latest_time = 0;
    
    // Time Complexity: O(N) - single pass
    for (int pos : ant_positions) {
        int distance_left = pos;
        int distance_right = pole_length - pos;
        
        int min_dist = min(distance_left, distance_right);
        int max_dist = max(distance_left, distance_right);
        if (min_dist > earliest_time) earliest_time = min_dist;
        if (max_dist > latest_time) latest_time = max_dist;
    }
    
    cout << earliest_time << " " << latest_time << "\n";
}

int main() {
    int num_test_cases;
    cin >> num_test_cases;
    while (num_test_cases--) {
        int pole_length, num_ants;
        cin >> pole_length >> num_ants;
        // Space Complexity: O(N)
        vector<int> ant_positions(num_ants);
        for (int i = 0; i < num_ants; ++i) {
            cin >> ant_positions[i];
        }
        
        calculateAntTimes(pole_length, ant_positions);
    }
    return 0;
}
\end{lstlisting}